\documentclass[12pt, a4paper]{article}

% --- Packages ---
\usepackage[utf8]{inputenc}
\usepackage[T1]{fontenc}
\usepackage{geometry}
\geometry{left=2.5cm, right=2.5cm, top=2.5cm, bottom=2.5cm}
\usepackage{times}
\usepackage{natbib}
\setcitestyle{authoryear,open={(},close={)}}
\usepackage{hyperref}
\usepackage{graphicx}
\usepackage{titlesec}
\usepackage{enumitem}
\usepackage{setspace}
\usepackage{amsmath}
\usepackage{amssymb}
\usepackage{booktabs}
\usepackage{array}
\usepackage{multirow}
\usepackage{xcolor}
\usepackage{mdframed}
\usepackage{soul}

\hypersetup{
    colorlinks=true,
    linkcolor=blue,
    filecolor=magenta,
    urlcolor=cyan,
    citecolor=blue,
    pdfencoding=auto
}

\onehalfspacing
\setlength{\parskip}{0.5em}
\titlespacing{\section}{0pt}{14pt}{6pt}
\titlespacing{\subsection}{0pt}{10pt}{4pt}
\titlespacing{\subsubsection}{0pt}{8pt}{3pt}

% Outline annotation environment
\newmdenv[
  leftmargin=0pt,
  rightmargin=0pt,
  innerleftmargin=8pt,
  innerrightmargin=8pt,
  innertopmargin=4pt,
  innerbottommargin=4pt,
  linecolor=gray!40,
  backgroundcolor=gray!6,
  linewidth=0.5pt
]{outline_note}

% =========================================================
\title{\textbf{Multi-Agent Reinforcement Learning for Autonomous Driving:\\
Algorithms, Coordination, and Emerging Paradigms}\\[0.5em]
{\large \textit{Survey Outline}}}
\author{\large Literature Review}
\date{\today}

\begin{document}

\maketitle

\begin{outline_note}
\textbf{Outline Purpose.}
This document presents the structural skeleton of the full survey.
Each section heading is accompanied by (1) a brief statement of its \emph{scope},
(2) the \emph{key papers} to be discussed, and (3) the \emph{planned figures or
tables} to be included.
Paper tags follow the format \texttt{[Ref\,N]} matching the 100-paper reference list.
\end{outline_note}

\tableofcontents
\newpage

% ==========================================================
\section*{Abstract \textnormal{\small (planned: $\approx$350 words)}}
\addcontentsline{toc}{section}{Abstract}

\begin{outline_note}
\textbf{Content.}
Provide a crisp problem statement (MARL for AD), explicitly state the survey
scope (algorithms, communication, simulation, applications), the literature
selection methodology (systematic search, 100 primary references spanning
2003--2024), and the three-paradigm organizing framework.  Explicitly enumerate
the \emph{three taxonomic layers}: (i) foundational single-agent RL; (ii) MARL
algorithms and coordination protocols; (iii) scalable and
Transformer/LLM-augmented agents.  Conclude with a sentence on the
survey's unique contribution versus existing reviews.
\end{outline_note}

% ==========================================================
\section{Introduction \textnormal{\small ($\approx$2{,}000 words)}}
\label{sec:intro}

\subsection{Motivation and Problem Statement}

\begin{outline_note}
\textbf{Scope.}
Define autonomous driving as a canonical multi-agent decision problem.
Describe the core difficulty: ego-vehicle actions and surrounding agents'
actions are mutually dependent, making single-agent POMDP solutions
insufficient.  Introduce the survey's central research question:
\emph{How have MARL algorithms, coordination protocols, and cognitive
architectures evolved to address the cooperative and competitive challenges
inherent in multi-agent traffic environments?}

\textbf{Key papers.}
\begin{itemize}[nosep]
  \item General MARL survey: \texttt{[Ref\,8]} (Zhang et al.\ 2021),
        \texttt{[Ref\,7]} (Busoniu et al.\ 2008)
  \item Autonomous driving RL survey: \texttt{[Ref\,12]} (Sallab et al.\ 2017);
        Kiran et al.\ 2021 (existing review)
  \item Driving-specific motivation: \texttt{[Ref\,56]} (Shalev-Shwartz et al.\ 2016)
\end{itemize}

\textbf{Planned figure.}
\textbf{Figure 1} -- A horizontal three-paradigm evolution timeline (2003--2024)
with color-coded swim lanes (blue: single-agent foundations, green: MARL \&
coordination, orange: Transformer/LLM era).  Each node represents a milestone
paper; curved arrows between lanes label the specific limitation that drove
the transition.
\end{outline_note}

\subsection{Driving Forces Behind Paradigm Shifts}

\begin{outline_note}
\textbf{Scope.}
Analyze \emph{why} the field evolved, not merely \emph{what} changed.
Identify four categories of drivers:
(1) algorithmic breakthroughs (deep RL, GNN, Transformer),
(2) hardware and compute scaling,
(3) dataset and simulation ecosystem maturation (\texttt{[Ref\,72--75,82]}),
(4) real-world deployment demands (Robotaxi fleets, V2X infrastructure).

\textbf{Planned table.}
\textbf{Table 1} -- Paradigm shift summary: paradigm name, time period,
core limitation of predecessor, driving innovation, representative benchmark.
\end{outline_note}

\subsection{Survey Scope and Literature Selection Methodology}

\begin{outline_note}
\textbf{Scope.}
Explicitly state:
\begin{itemize}[nosep]
  \item \textbf{Topic coverage}: MARL algorithms (value decomposition, actor-critic,
        communication), trajectory prediction, multi-vehicle cooperative control,
        traffic management, sim-to-real transfer, safety, and LLM-based agents.
  \item \textbf{Time span}: 1993--2024 (foundational theory to present).
  \item \textbf{Search strategy}: IEEE Xplore, ACM DL, arXiv; keywords
        \emph{``multi-agent reinforcement learning autonomous driving''},
        \emph{``cooperative driving deep RL''}, \emph{``vehicle coordination MARL''};
        backward/forward citation tracing.
  \item \textbf{Inclusion criteria}: peer-reviewed or highly cited preprints;
        direct relevance to either MARL methodology or AD application.
  \item \textbf{Exclusion criteria}: purely perception-only methods, rule-based
        planners without learning, non-vehicular robotics.
  \item Total: 100 primary references.
\end{itemize}

\textbf{Planned figure.}
\textbf{Figure 2} -- PRISMA-style flowchart showing the literature selection
process (initial hits, de-duplication, relevance screening, final 100).
\end{outline_note}

\subsection{Three-Layer Taxonomy}

\begin{outline_note}
\textbf{Scope.}
Present the organizing framework of the survey:
\begin{enumerate}[nosep]
  \item \textbf{Layer 1 -- Algorithmic Foundations}: game-theoretic formalisms,
        value-based MARL, policy-gradient MARL, hierarchical and role-based methods.
  \item \textbf{Layer 2 -- Coordination and Communication}: CTDE vs.\ decentralized
        execution, communication-learning, V2X cooperative perception.
  \item \textbf{Layer 3 -- Scalable and Cognitive Architectures}:
        Transformer-based MARL, LLM-driven cognitive agents, sim-to-real transfer.
\end{enumerate}
Emphasize how the three layers map onto the three paradigms in
Subsection~1.1.

\textbf{Planned figure.}
\textbf{Figure 3} -- Hierarchical taxonomy tree: root node
``MARL for Autonomous Driving'', three main branches, leaf nodes labeled with
representative method names and publication years.
\end{outline_note}

\subsection{Comparison with Existing Surveys and Contributions}

\begin{outline_note}
\textbf{Scope.}
Contrast this survey with:
(a) general MARL surveys (\texttt{[Ref\,7,8]}) -- lack AD-specific coverage;
(b) AD-RL surveys (Kiran et al.\ 2021, \texttt{[Ref\,12]}) -- primarily
single-agent, limited MARL depth;
(c) cooperative driving reviews -- narrow scope (often only traffic signal or only
V2X).
State the four unique contributions: (1) unified three-layer taxonomy,
(2) systematic comparison of 100 papers, (3) critical limitation analysis per method
category, (4) structured open-challenges vs.\ future-directions mapping.
\end{outline_note}

\subsection{Paper Organization}

\begin{outline_note}
One paragraph describing the section flow with explicit cross-references.
\end{outline_note}

% ==========================================================
\section{Background and Formal Foundations \textnormal{\small ($\approx$2{,}500 words)}}
\label{sec:background}

\subsection{Single-Agent Formalism: MDP and POMDP}

\begin{outline_note}
\textbf{Scope.}
Define MDP tuple $\langle \mathcal{S}, \mathcal{A}, P, R, \gamma \rangle$.
Motivate the POMDP extension for driving (limited sensor range, occluded
agents).  Define belief-state planning and its computational intractability.

\textbf{Key papers.}
\texttt{[Ref\,70]} (Hubmann et al.\ 2018 -- POMDP autonomous driving).
\end{outline_note}

\subsection{Multi-Agent Extension: Stochastic Games and Dec-POMDP}

\begin{outline_note}
\textbf{Scope.}
Define Stochastic (Markov) Game \texttt{[Ref\,11]} (Shapley 1953),
Markov Game \texttt{[Ref\,10]} (Littman 1994).
Define Dec-POMDP \texttt{[Ref\,9]} (Bernstein et al.\ 2002): joint action space,
individual observations, shared or individual rewards.
Discuss NEXP-hardness and practical relaxations.
Show how cooperative AD maps to Dec-POMDP and mixed-motive AD maps to
Markov Games.

\textbf{Planned figure.}
\textbf{Figure 4} -- Side-by-side comparison of POMDP (single agent) and
Dec-POMDP/POSG (multi-agent): observation model, reward structure,
execution policy structure.
\end{outline_note}

\subsection{Deep Reinforcement Learning Primer}

\begin{outline_note}
\textbf{Scope.}
Briefly survey the algorithmic building blocks used throughout the paper:
Q-learning (DQN), policy gradients (REINFORCE), Actor-Critic (A3C/A2C),
PPO \texttt{[Ref\,Schulman 2017]}, DDPG/SAC for continuous control.
Emphasize the aspects relevant to multi-agent extensions.
\end{outline_note}

\subsection{Graph Neural Networks for Relational Reasoning}

\begin{outline_note}
\textbf{Scope.}
Motivate GNNs for traffic interaction modeling.
Cover Graph Attention Networks \texttt{[Ref\,17]} (Velickovic et al.\ 2018),
their role in encoding spatial agent relationships.
This subsection sets up the representation methods in Section~4.

\textbf{Key papers.}
\texttt{[Ref\,17]} (GAT), \texttt{[Ref\,31]} (Graph Conv RL),
\texttt{[Ref\,16]} (Trajectron++).
\end{outline_note}

\subsection{Overview of Autonomous Driving Software Stack}

\begin{outline_note}
\textbf{Scope.}
Describe modular pipeline (perception $\to$ prediction $\to$ planning $\to$
control) and end-to-end alternatives.
Clarify where MARL intervenes in this stack (prediction, planning,
or full end-to-end).

\textbf{Key papers.}
\texttt{[Ref\,13]} (Codevilla CILR 2018 -- conditional imitation learning),
\texttt{[Ref\,12]} (Sallab 2017 -- deep RL for AD).
\end{outline_note}

% ==========================================================
\section{Paradigm I: Single-Agent Learning Foundations \textnormal{\small ($\approx$2{,}000 words)}}
\label{sec:paradigm1}

\begin{outline_note}
\textbf{Paradigm framing.}
This paradigm (approx.\ 2016--2018) establishes the core learning machinery
later extended to multi-agent settings.  Critically discuss its central
limitation: single-agent learning ignores the non-stationarity arising
from other learning agents, making it unsuitable as a standalone solution
for dense traffic.
\end{outline_note}

\subsection{Imitation Learning and Behavior Cloning}

\subsubsection{End-to-End Learning: DAVE-2 and Successors}

\begin{outline_note}
\textbf{Key papers.}
DAVE-2 / NVIDIA end-to-end \texttt{(Bojarski et al.\ 2016 -- in refs folder)},
\texttt{[Ref\,13]} (Codevilla 2018).
\textbf{Contributions}: raw pixel to steering command; data efficiency.
\textbf{Limitations}: covariate shift; brittle recovery; no explicit
agent modeling.
\end{outline_note}

\subsubsection{Addressing Covariate Shift: DAgger and Data Augmentation}

\begin{outline_note}
\textbf{Key papers.}
ChauffeurNet \texttt{(Bansal et al.\ 2018 -- in refs folder)},
DAgger (Ross et al.\ 2011).
\textbf{Contributions}: perturbation-based synthesis; demonstrates that
30M expert demonstrations are insufficient without augmentation.
\textbf{Limitations}: still single-agent; synthetic perturbations
do not capture realistic multi-agent interaction.
\end{outline_note}

\subsection{Reinforcement Learning for Single-Agent Driving}

\begin{outline_note}
\textbf{Key papers.}
\texttt{[Ref\,12]} (Sallab 2017), \texttt{[Ref\,71]} (Kahn risk-sensitive 2017),
\texttt{[Ref\,60]} (Chen hierarchical ICRA 2019),
PPO \texttt{(Schulman 2017 -- in refs folder)}.
\textbf{Coverage}: model-free on-policy (PPO), off-policy (DDPG),
risk-sensitive RL, hierarchical RL for lane change and intersection.
\textbf{Limitations}: non-stationarity when other vehicles are also
learning; scalability to dense scenarios.
\end{outline_note}

\subsection{Interaction-Aware Trajectory Prediction (Single-Agent Perspective)}

\begin{outline_note}
\textbf{Key papers.}
Social LSTM \texttt{[Ref\,14]} (Alahi 2016),
DESIRE \texttt{[Ref\,15]} (Lee 2017),
Social GAN \texttt{[Ref\,61]} (Gupta 2018).
\textbf{Coverage}: grid social pooling, GAN-based multi-modal prediction.
\textbf{Limitations}: prediction is decoupled from planning; ignores
the ego-agent's own reactive behavior in predictions.

\textbf{Planned figure.}
\textbf{Figure 5} -- Social LSTM grid pooling versus GAN-based stochastic
trajectory sampling: architectural comparison diagram.
\end{outline_note}

\subsection{Critical Analysis: Limitations Driving Paradigm~II}

\begin{outline_note}
Summarize three concrete failure modes:
(1) covariate shift in IL, (2) non-stationarity in RL, (3) prediction-planning
disconnect.  Each limitation maps directly to a solution introduced in
Paradigm II, forming the causal bridge between sections.
\end{outline_note}

% ==========================================================
\section{Paradigm II: Multi-Agent Reinforcement Learning -- Algorithms and Coordination \textnormal{\small ($\approx$6{,}000 words)}}
\label{sec:paradigm2}

\begin{outline_note}
\textbf{Paradigm framing.}
This is the core of the survey (approx.\ 2017--2022).
It is organized into four thematic subsections: (A) MARL algorithms,
(B) communication and coordination, (C) interaction-aware representation
and prediction-planning integration, and (D) autonomous driving applications.
\end{outline_note}

% ---- 4.A ----
\subsection{Multi-Agent RL Algorithms}
\label{sec:marl_alg}

\subsubsection{Independent Learning and Its Limitations}

\begin{outline_note}
\textbf{Key papers.}
Independent Q-Learning (IQL) \texttt{[Ref\,5]} (Tan 1993),
Independent PPO (IPPO) \texttt{[Ref\,42]}.
\textbf{Contributions}: simplicity, parameter sharing, scalable.
\textbf{Limitations}: environment non-stationarity from each agent's
perspective; convergence guarantees lost.
\end{outline_note}

\subsubsection{CTDE: Centralized Training, Decentralized Execution}

\begin{outline_note}
\textbf{Framing.}
Introduce the CTDE paradigm \texttt{[Ref\,6]} (Kraemer \& Banerjee 2016) as
the dominant design philosophy: a centralized critic with access to global
state at training time; each actor uses only local observations at
execution.

\textbf{Value Decomposition Methods.}
\begin{itemize}[nosep]
  \item VDN \texttt{[Ref\,3]} (Sunehag 2018): additive decomposition;
        credit assignment via gradient.
        \textbf{Limitation}: linear decomposition too restrictive.
  \item QMIX \texttt{[Ref\,2]} (Rashid 2018): monotonic mixing network;
        IGM condition.  \textbf{Limitation}: monotonic constraint
        cannot represent all joint Q-functions.
  \item QTRAN \texttt{[Ref\,19]} (Son 2019): relaxes monotonicity.
        \textbf{Limitation}: optimization instability.
  \item MAVEN \texttt{[Ref\,21]} (Mahajan 2019): latent variable
        exploration.
\end{itemize}

\textbf{Actor-Critic CTDE Methods.}
\begin{itemize}[nosep]
  \item MADDPG \texttt{[Ref\,1,24]} (Lowe 2017): mixed cooperative-competitive;
        centralized critic per agent.
  \item COMA \texttt{[Ref\,4,36]} (Foerster 2018): counterfactual baseline
        for credit assignment.
  \item MAPPO \texttt{[Ref\,18]} (Yu 2021): PPO in MARL; surprisingly strong
        baseline.  \textbf{Limitation}: on-policy; sample inefficiency.
  \item HATRPO \texttt{[Ref\,23]} (Kuba 2021): trust-region extension.
  \item FACMAC \texttt{[Ref\,40]} (Peng 2021): continuous action
        decomposition.
\end{itemize}

\textbf{Planned table.}
\textbf{Table 2} -- CTDE algorithm comparison: algorithm, decomposition type,
action space (discrete/continuous), credit assignment mechanism,
key limitation, primary benchmark.
\end{outline_note}

\subsubsection{Role-Based and Structured MARL}

\begin{outline_note}
\textbf{Key papers.}
RODE \texttt{[Ref\,20]} (Wang 2021): action-space role decomposition.
ROMA \texttt{[Ref\,22]} (Wang 2020): role regularization via mutual information.
\textbf{Relevance to AD}: heterogeneous agents (trucks, motorcycles, pedestrians)
naturally map to distinct roles.
\textbf{Limitations}: role discovery is task-sensitive; hard to generalize
across traffic scenarios.
\end{outline_note}

\subsubsection{Hierarchical MARL}

\begin{outline_note}
\textbf{Key papers.}
FeUdal Networks \texttt{[Ref\,47]} (Vezhnevets 2017),
Hierarchical MARL for traffic \texttt{[Ref\,48]} (Chu 2020),
hierarchical RL for urban driving \texttt{[Ref\,60]}.
\textbf{Framing}: high-level goal selection decoupled from low-level
primitive execution enables temporal abstraction.
\textbf{Limitations}: reward decomposition across levels is non-trivial;
sub-goal definitions require domain knowledge.
\end{outline_note}

\subsubsection{Game-Theoretic Approaches}

\begin{outline_note}
\textbf{Key papers.}
Stochastic games \texttt{[Ref\,11]},
Markov games \texttt{[Ref\,10]},
Nash Q-learning \texttt{[Ref\,49]} (Hu \& Wellman 2003),
LOLA \texttt{[Ref\,50]} (Foerster 2018),
PSRO \texttt{[Ref\,45]} (Lanctot 2017).
\textbf{Relevance to AD}: mixed-motive settings (e.g., merging, yielding)
require Nash equilibrium reasoning.
\textbf{Limitations}: Nash computation is intractable at scale; LOLA
requires differentiating through opponents' update steps.
\end{outline_note}

\subsubsection{Population-Based Training and Self-Play}

\begin{outline_note}
\textbf{Key papers.}
AlphaStar \texttt{[Ref\,46]}, PSRO \texttt{[Ref\,45]}.
\textbf{Relevance to AD}: generating diverse, robust opponents
for stress-testing driving agents.
\textbf{Limitations}: high compute cost; opponent distribution may
not match real road users.
\end{outline_note}

% ---- 4.B ----
\subsection{Communication and Coordination Protocols}
\label{sec:comm}

\subsubsection{Learned Communication: Early Approaches}

\begin{outline_note}
\textbf{Key papers.}
DIAL/RIAL \texttt{[Ref\,26]} (Foerster 2016),
CommNet \texttt{[Ref\,27]} (Sukhbaatar 2016).
\textbf{Contributions}: first demonstration that communication can be
jointly learned with policy.
\textbf{Limitations}: no bandwidth constraint; broadcast communication
does not scale.
\end{outline_note}

\subsubsection{Targeted and Selective Communication}

\begin{outline_note}
\textbf{Key papers.}
TarMAC \texttt{[Ref\,28]} (Das 2019): attention-based targeted messages.
ATOC \texttt{[Ref\,29]} (Jiang 2018): learned gating of communication.
VBC \texttt{[Ref\,34]} (Wang 2021): value-based selection.
IC3Net \texttt{[Ref\,35]} (Singh 2019): independently controlled gating.
\textbf{Limitations}: selective communication introduces non-differentiability;
message encoding/decoding can become an information bottleneck.
\end{outline_note}

\subsubsection{Graph-Based Communication}

\begin{outline_note}
\textbf{Key papers.}
Graph Conv RL \texttt{[Ref\,31]} (Jiang 2019): GCN over agent interaction graph.
Structured communication \texttt{[Ref\,30]} (Kim 2019): learned discrete graph.
\textbf{Contributions}: scalable, topology-aware message passing.
\textbf{Limitations}: dynamic graph topology changes as agents enter/leave
communication range.

\textbf{Planned figure.}
\textbf{Figure 6} -- Communication topology evolution: broadcast
$\to$ attention-gated $\to$ graph-structured.  Diagram showing
message flow for a 5-vehicle scenario under each paradigm.
\end{outline_note}

\subsubsection{Bandwidth-Constrained Communication}

\begin{outline_note}
\textbf{Key papers.}
Message-Dropout \texttt{[Ref\,32]} (Kim 2019),
limited-bandwidth MARL \texttt{[Ref\,33]} (Wang 2020).
\textbf{Relevance to AD}: V2X channel capacity is finite; packet loss
is realistic in urban canyons.
\textbf{Limitations}: current methods simulate packet loss but do not
model realistic wireless channel models.
\end{outline_note}

\subsubsection{V2X Cooperative Perception}

\begin{outline_note}
\textbf{Key papers.}
V2VNet \texttt{(Wang et al.\ 2020 -- in refs folder)}: intermediate feature
fusion via compressed spatial maps.
\textbf{Contributions}: V2V feature sharing improves object detection
at long range.
\textbf{Limitations}: feature representations must be aligned across
heterogeneous sensor suites; requires time-synchronized maps.

\textbf{Planned figure.}
\textbf{Figure 7} -- V2X fusion strategies: early fusion (raw sensor data),
intermediate fusion (feature maps), late fusion (predictions).
Comparison of detection range and communication bandwidth for each.
\end{outline_note}

% ---- 4.C ----
\subsection{Interaction-Aware Representation and Prediction-Planning Integration}
\label{sec:representation}

\subsubsection{GNN-Based Trajectory Prediction}

\begin{outline_note}
\textbf{Key papers.}
Trajectron++ \texttt{[Ref\,16]} (Salzmann 2020): heterogeneous agent types,
temporal GNN.
VectorNet \texttt{[Ref\,62]} (Gao 2020): vectorized HD-map + agent trajectory
encoding.
LaneGCN \texttt{[Ref\,63]} (Liang 2020): lane-graph convolution.
TNT \texttt{[Ref\,65]} (Zhao 2020): target-conditioned trajectory sampling.
MultiPath \texttt{[Ref\,64]} (Chai 2019): anchor-based multi-modal prediction.
\textbf{Limitations}: these models predict each agent's trajectory
\emph{independently} conditioned on others' past motion, ignoring
the coupled nature of future decisions.

\textbf{Planned table.}
\textbf{Table 3} -- Trajectory prediction methods: input representation,
interaction model, multi-modality mechanism, benchmark, ADE/FDE.
\end{outline_note}

\subsubsection{Prediction-Planning Bidirectional Coupling}

\begin{outline_note}
\textbf{Key papers.}
PiP \texttt{(Song et al.\ 2020 -- in refs folder)}: planning-informed
prediction feeds candidate ego plans into the prediction network.
\textbf{Contribution}: breaks the one-directional prediction$\to$planning
pipeline; ego plan affects surrounding agent trajectory estimates.
\textbf{Limitation}: closed-loop interaction between planning and prediction
is approximated rather than exact.
\end{outline_note}

\subsubsection{Game-Theoretic Motion Planning}

\begin{outline_note}
\textbf{Key papers.}
Game-theoretic motion planning \texttt{[Ref\,94]} (Fridovich-Keil 2020).
Game-theoretic autonomous driving \texttt{[Ref\,58]} (Fisac 2019).
Interactive decision making via deep RL \texttt{[Ref\,59]} (Sadigh 2016).
Intent-aware planning \texttt{[Ref\,69]} (Schwarting 2021).
\textbf{Coverage}: level-$k$ reasoning, potential-game formulations,
responsibility-sensitive safety (RSS).
\textbf{Limitations}: game-theoretic solvers require known utility functions
and agent rationality assumptions that real humans violate.

\textbf{Planned figure.}
\textbf{Figure 8} -- Level-$k$ game-theoretic reasoning at a merge:
level-0 agent acts selfishly; level-1 anticipates level-0; level-2
anticipates level-1.
\end{outline_note}

% ---- 4.D ----
\subsection{Autonomous Driving Applications of MARL}
\label{sec:ad_apps}

\subsubsection{Traffic Signal Control}

\begin{outline_note}
\textbf{Key papers.}
CoLight \texttt{[Ref\,51]} (Wei 2019): graph attention for network-level
signal coordination.
PressLight \texttt{[Ref\,52]} (Wei 2019): max-pressure theory + RL.
FRAP \texttt{[Ref\,53]} (Zheng 2019): factorized phase representation.
Large-scale hierarchical MARL \texttt{[Ref\,54]} (Chu 2020).
\textbf{Comparison}: rule-based pressure control vs.\ learned coordination;
scalability vs.\ optimality trade-off.
\textbf{Limitations}: ignores vehicle heterogeneity (emergency vehicles,
public transit); reward design with network externalities is unsolved.

\textbf{Planned table.}
\textbf{Table 4} -- Traffic signal MARL methods: number of intersections,
state/action space, coordination mechanism, improvement over fixed-time
baseline.
\end{outline_note}

\subsubsection{Multi-Vehicle Cooperative Driving}

\begin{outline_note}
\textbf{Key papers.}
Cooperative lane change \texttt{[Ref\,55]} (Wang 2018 T-ITS).
Safe multi-agent RL \texttt{[Ref\,56]} (Shalev-Shwartz 2016).
Learning interactive driving policies \texttt{[Ref\,57]} (Schwarting 2019).
\textbf{Coverage}: highway lane change, on-ramp merging, cooperative
intersection crossing.
\textbf{Limitations}: methods assume perfect V2V communication; real-world
latency and packet loss degrade coordination performance.
\end{outline_note}

\subsubsection{Mixed-Autonomy Traffic}

\begin{outline_note}
\textbf{Key papers.}
Flow \texttt{[Ref\,66]} (Wu 2017): RL for mixed human-AV traffic control.
Stabilizing traffic with AVs \texttt{[Ref\,67]} (Wu 2018 Nature).
Autonomous driving in mixed-autonomy systems \texttt{[Ref\,100]}
(Talebpour 2016).
\textbf{Contribution}: even a small fraction of AVs can suppress
stop-and-go waves when coordinated.
\textbf{Limitations}: assumes AVs can communicate; human drivers are
modeled by simple car-following models that underestimate behavioral
diversity.
\end{outline_note}

\subsubsection{Highway and Urban Driving with MARL}

\begin{outline_note}
\textbf{Key papers.}
MARL for highway driving \texttt{[Ref\,68]} (Isele 2018).
POMDP-based autonomous driving \texttt{[Ref\,70]} (Hubmann 2018).
Multi-agent connected AV \texttt{(Palanisamy 2020 -- in refs folder)}.
\textbf{Coverage}: lane selection, speed negotiation, unsignalized
intersection management.
\textbf{Limitations}: open-world generalization; long-tail scenarios
(unexpected pedestrian crossings) remain unsolved.
\end{outline_note}

% ---- 4.E ----
\subsection{Simulation Platforms and Benchmarks}
\label{sec:sim}

\begin{outline_note}
\textbf{Key papers.}
CARLA \texttt{[Ref\,72]} (Dosovitskiy 2017 -- also in refs folder).
SUMO \texttt{[Ref\,73]} (Krajzewicz 2012).
SMARTS \texttt{[Ref\,82]} / \texttt{(Zhou 2020 -- in refs folder)}.
Waymo Open Motion Dataset \texttt{[Ref\,74]} (Ettinger 2021).
nuScenes \texttt{[Ref\,75]} / NuPlan \texttt{(Caesar 2022 -- in refs folder)}.
Argoverse \texttt{[Ref\,97]} (Chang 2019).
SMAC \texttt{[Ref\,82]} (Samvelyan 2019).

\textbf{Planned table.}
\textbf{Table 5} -- Simulation and benchmark ecosystem:
platform name, year, sensor modalities, agent types, MARL support,
primary task, open-source availability.
\end{outline_note}

\subsection{Critical Analysis: Limitations Driving Paradigm~III}

\begin{outline_note}
Identify three concrete failure modes of Paradigm II:
(1) sample inefficiency and training instability at scale (non-stationarity
persists even under CTDE);
(2) long-tail generalization failure (MARL agents overfit to training
scenarios, failing in rare but safety-critical edge cases);
(3) interpretability deficit (complex MARL policies are opaque;
regulatory acceptance requires explainability).
\end{outline_note}

% ==========================================================
\section{Paradigm III: Transformer-Based and LLM-Augmented MARL \textnormal{\small ($\approx$3{,}000 words)}}
\label{sec:paradigm3}

\begin{outline_note}
\textbf{Paradigm framing.}
This paradigm (approx.\ 2021--present) addresses interpretability and
long-tail generalization by importing architectural innovations from NLP
and by exploiting the commonsense knowledge encoded in large language models.
\end{outline_note}

\subsection{Transformer Architectures for MARL}

\subsubsection{Sequential Decision Making as Sequence Modeling}

\begin{outline_note}
\textbf{Key papers.}
Decision Transformer \texttt{[Ref\,78]} (Chen 2021): offline RL as
autoregressive sequence prediction.
Trajectory Transformer \texttt{[Ref\,79]} (Janner 2021): full
trajectory-level distribution modeling.
GTrXL \texttt{[Ref\,77]} (Parisotto 2020): stabilized Transformer for RL.
\textbf{Contributions}: eliminates explicit reward function design;
enables offline learning from diverse driving logs.
\textbf{Limitations}: purely offline; requires return-labeled datasets;
sub-optimal for online fine-tuning.
\end{outline_note}

\subsubsection{Multi-Agent Transformer (MAT)}

\begin{outline_note}
\textbf{Key papers.}
MAT \texttt{[Ref\,76]} (Wen 2022): encoder-decoder Transformer modeling
sequential joint-action generation; global dependency via attention.
R-MAT \texttt{[Ref\,80]} (Papoudakis 2021): scalable attention structure.
\textbf{Contributions}: captures global agent dependencies without
hand-crafted communication graphs.
\textbf{Limitations}: quadratic attention complexity with agent count;
not yet demonstrated in high-fidelity AD simulators at scale.

\textbf{Planned figure.}
\textbf{Figure 9} -- MAT encoder-decoder architecture: agent
observations encoded jointly; decoder autoregressively generates
actions conditioned on previous agents' actions.
\end{outline_note}

\subsubsection{GNN + Transformer Hybrid Models}

\begin{outline_note}
\textbf{Key papers.}
Deep Graph Networks for MARL \texttt{[Ref\,81]} (Li 2020).
\textbf{Framing}: GNNs model local topology; Transformers model
global context; hybrids combine both.
\textbf{Limitations}: increased model complexity; requires careful
positional encoding for non-Euclidean graph structure.
\end{outline_note}

\subsection{Large Language Models as Cognitive Driving Agents}

\subsubsection{LLM as Commonsense Reasoner}

\begin{outline_note}
\textbf{Key papers.}
Drive Like a Human \texttt{(Fu et al.\ 2023 -- in refs folder)}:
GPT-4 zero-shot chain-of-thought reasoning for novel scenarios.
Voyager \texttt{[Ref\,88]} (Wang 2023): skill-accumulation agent.
\textbf{Contributions}: LLM knowledge covers long-tail situations absent
from training datasets; reasoning trace provides interpretability.
\textbf{Limitations}: hallucination risk; latency ($>$100\,ms per query)
incompatible with real-time control loops ($<$50\,ms).
\end{outline_note}

\subsubsection{Memory-Augmented Lifelong Adaptation}

\begin{outline_note}
\textbf{Key papers.}
DiLu \texttt{(Wen et al.\ 2024 -- in refs folder)}: memory module
stores past driving episodes; reflection mechanism updates policy.
\textbf{Contributions}: bridges the gap between static pre-training
and continuous deployment; reduces catastrophic forgetting.
\textbf{Limitations}: memory retrieval quality degrades with corpus size;
episode similarity metrics require calibration.
\end{outline_note}

\subsubsection{LLM as Motion Planner}

\begin{outline_note}
\textbf{Key papers.}
GPT-Driver \texttt{(Mao et al.\ 2023 -- in refs folder)}:
waypoint generation framed as language completion.
LLM-Planner \texttt{[Ref\,90]} (Song 2023): high-level route planning.
\textbf{Contributions}: leverages LLM structured output parsing for
interpretable intermediate waypoints.
\textbf{Limitations}: language-to-coordinate mapping is imprecise;
no safety guarantee on generated trajectories.
\end{outline_note}

\subsubsection{Neuro-Symbolic Hybrid: LLM + Model Predictive Control}

\begin{outline_note}
\textbf{Key papers.}
LanguageMPC \texttt{(Sha et al.\ 2025 -- in refs folder)}: LLM generates
behavioral parameters consumed by a downstream MPC controller.
\textbf{Contributions}: separates semantic reasoning (LLM) from
real-time trajectory optimization (MPC); provides formal safety guarantees
at the execution level.
\textbf{Limitations}: MPC formulation is hand-crafted per scenario;
the LLM $\to$ MPC interface requires extensive prompt engineering.

\textbf{Planned figure.}
\textbf{Figure 10} -- Neuro-symbolic architecture: LLM reasoning layer
generates cost weights; MPC execution layer solves constrained optimization
with hard safety constraints.
\end{outline_note}

\subsection{Safety and Robustness in MARL}
\label{sec:safety}

\begin{outline_note}
\textbf{Key papers.}
CPO \texttt{[Ref\,83]} (Achiam 2017): constrained policy optimization.
Safe RL via shielding \texttt{[Ref\,84]} (Alshiekh 2018): formal safety filter.
Risk-aware MARL \texttt{[Ref\,85]} (Prashanth 2016): CVaR optimization.
\textbf{Coverage}: constraint formulation (CMDP), shielding with
temporal logic specifications, risk-sensitive objectives.
\textbf{Limitations}: formal verification does not scale to neural policies;
CVaR optimization is conservative; online shielding adds latency.

\textbf{Planned table.}
\textbf{Table 6} -- Safety mechanisms: approach, formal guarantee,
computational overhead, coverage of stochastic failures.
\end{outline_note}

\subsection{Sim-to-Real Transfer and Domain Generalization}

\begin{outline_note}
\textbf{Key papers.}
Domain randomization \texttt{[Ref\,86]} (Tobin 2017).
MAML meta-RL \texttt{[Ref\,87]} (Finn 2017).
\textbf{Coverage}: randomization of lighting, weather, traffic density;
meta-learning for rapid adaptation to new environments.
\textbf{Limitations}: the sim-to-real gap for complex social interactions
is larger than for perception; behavioral fidelity of simulated agents
remains poor.
\end{outline_note}

\subsection{Scalability and Generalization Across Agent Counts}

\begin{outline_note}
\textbf{Key papers.}
Deep Graph Networks \texttt{[Ref\,81]},
R-MAT \texttt{[Ref\,80]},
Interaction-aware prediction via GNN \texttt{[Ref\,96]} (Deo 2020).
\textbf{Coverage}: architectures that gracefully handle variable numbers
of agents (e.g., sparse highway vs.\ dense urban intersection).
\textbf{Limitations}: training on variable-agent environments requires
careful curriculum design.
\end{outline_note}

% ==========================================================
\section{Cross-Cutting Analysis \textnormal{\small ($\approx$2{,}000 words)}}
\label{sec:crosscut}

\subsection{Unified Comparison of MARL Approaches}

\begin{outline_note}
\textbf{Planned table.}
\textbf{Table 7} -- Master comparison table covering all three paradigms:
method, year, paradigm layer, algorithm class, AD scenario, benchmark,
key metric, communication required, safety mechanism, primary limitation.
This table spans all 100 references.
\end{outline_note}

\subsection{Credit Assignment: From Tabular to Neural}

\begin{outline_note}
\textbf{Key papers.}
VDN, QMIX, COMA \texttt{[Ref\,2,3,4]},
Counterfactual credit assignment \texttt{[Ref\,36]},
Shapley value MARL \texttt{[Ref\,91]} (Wang 2020),
Lazy agents problem \texttt{[Ref\,37]}.
\textbf{Planned figure.}
\textbf{Figure 11} -- Credit assignment taxonomy: global reward
$\to$ difference rewards $\to$ counterfactual baselines
$\to$ Shapley values.  Complexity vs.\ accuracy trade-off curve.
\end{outline_note}

\subsection{Non-Stationarity: Diagnosis and Mitigation}

\begin{outline_note}
Systematic treatment of non-stationarity under:
(a) independent learning (convergence failure),
(b) CTDE (mitigated but not eliminated by centralized critic),
(c) opponent modeling \texttt{[Ref\,93]} (He 2016),
(d) learning with imagined opponent \texttt{[Ref\,50]}.
\end{outline_note}

\subsection{Reward Shaping for Traffic}

\begin{outline_note}
Survey reward functions used across AD applications:
speed, comfort, collision penalty, throughput, fairness (social welfare).
Discuss reward hacking examples and mitigation strategies.
\end{outline_note}

% ==========================================================
\section{Open Challenges \textnormal{\small ($\approx$1{,}500 words)}}
\label{sec:challenges}

\begin{outline_note}
\textbf{Note on structure.}
Each challenge is stated as a \emph{current unsolved problem};
the proposed research directions are deferred to Section~8.
This two-section structure directly addresses the reviewer feedback that
challenges and directions were previously conflated.
\end{outline_note}

\subsection{Non-Stationarity and Scalability in Large Fleets}

\begin{outline_note}
As the number of agents increases beyond toy scenarios (e.g.,
$>$50 vehicles in a city-scale simulation), CTDE's centralized critic
becomes computationally infeasible.  IQL resorts to
non-stationary, unstable training.  No current method provides
both scalability and convergence guarantees simultaneously.

\textbf{Evidence}: \texttt{[Ref\,7,8,99]}, SMARTS experiments.
\end{outline_note}

\subsection{Credit Assignment in Sparse Reward Settings}

\begin{outline_note}
In cooperative driving tasks (e.g., reducing network congestion),
individual contributions to the team reward are difficult to isolate.
Counterfactual methods (\texttt{[Ref\,4]}) mitigate but do not eliminate
the credit assignment problem.  Shapley-based methods (\texttt{[Ref\,91]})
are exact but exponentially expensive.
\end{outline_note}

\subsection{Formal Safety Verification}

\begin{outline_note}
Neural MARL policies lack formal safety certificates.
Shielding approaches (\texttt{[Ref\,84]}) rely on hand-crafted
formal models that cannot capture the full state space of real traffic.
Distributional shift at test time can violate safety constraints
that held during training (\texttt{[Ref\,83,85]}).
\end{outline_note}

\subsection{Bandwidth-Constrained and Intermittent Communication}

\begin{outline_note}
Current communication-learning MARL assumes reliable, low-latency
channels (\texttt{[Ref\,26--35]}).  Real V2X deployments experience
packet loss, variable latency (5--200\,ms), and partial connectivity.
No current method is trained under a realistic wireless channel model.
\end{outline_note}

\subsection{Long-Tail Scenario Coverage and Distributional Robustness}

\begin{outline_note}
MARL policies overfit to the scenario distribution of the training
simulator.  CARLA/SUMO (\texttt{[Ref\,72,73]}) cover common scenarios
but underrepresent rare-but-critical events (debris on road, aggressive
pedestrian behavior, sensor spoofing).
\end{outline_note}

\subsection{Real-Time Inference for LLM-Based Agents}

\begin{outline_note}
GPT-4-class models require $\sim$500\,ms--2\,s per query
(\texttt{[Ref\,88,89,90]}).  Control systems require $<$50\,ms
response latency.  Efficient approximations (speculative decoding,
caching) have not been validated in closed-loop driving.
\end{outline_note}

\subsection{Human-Agent Interaction and Mixed Autonomy Behavior}

\begin{outline_note}
Human drivers do not follow rational game-theoretic models
(\texttt{[Ref\,57,59,100]}).  Current MARL methods model humans as
either fully rational or as fixed stochastic policies; neither
captures the full behavioral distribution observed on real roads.
\end{outline_note}

% ==========================================================
\section{Future Research Directions \textnormal{\small ($\approx$1{,}500 words)}}
\label{sec:future}

\begin{outline_note}
\textbf{Note on structure.}
Each direction is a concrete technical path mapped to one or more
open challenges from Section~7.  This pairing makes the link between
problem and proposed solution explicit.
\end{outline_note}

\subsection{Generative World Models for Long-Tail Simulation}

\begin{outline_note}
\textbf{Challenge addressed}: \S7.5 (long-tail scenarios).
Generative models (diffusion models, video prediction networks)
conditioned on rare-event descriptions could synthesize realistic
adversarial traffic scenarios for curriculum training.
\end{outline_note}

\subsection{Scalable MARL via Mean-Field Approximation and Attention}

\begin{outline_note}
\textbf{Challenge addressed}: \S7.1 (scalability).
Mean-field MARL approximates the joint policy as a product of
individual policies interacting via a mean-field term.  Combined
with attention mechanisms (\texttt{[Ref\,76,80]}), this could scale
to city-level fleets without centralized critics.
\end{outline_note}

\subsection{Formal-Methods-Augmented MARL}

\begin{outline_note}
\textbf{Challenge addressed}: \S7.3 (safety verification).
Integrating STL/LTL specifications as differentiable constraints
within the RL training loop could provide formal guarantees while
retaining learned flexibility.  Shield learning (not just shield
inference) is a promising direction.
\end{outline_note}

\subsection{Causal Reasoning and Interpretable Credit Assignment}

\begin{outline_note}
\textbf{Challenge addressed}: \S7.2 (credit assignment).
Causal influence diagrams could disentangle each agent's causal
contribution to collective outcomes, enabling interpretable
and more efficient credit assignment beyond counterfactual baselines.
\end{outline_note}

\subsection{Efficient LLM Deployment for Real-Time Control}

\begin{outline_note}
\textbf{Challenge addressed}: \S7.6 (LLM latency).
Small language models (SLMs) fine-tuned on driving data,
combined with speculative decoding, could bring LLM reasoning
latency into the feasible real-time range.
The LanguageMPC hybrid architecture (\texttt{[Ref in refs folder]}) is a
promising template: LLM provides offline-amortized policy parameters;
MPC provides real-time trajectory guarantees.
\end{outline_note}

\subsection{Federated Multi-Agent Learning for Privacy-Preserving V2X}

\begin{outline_note}
\textbf{Challenge addressed}: \S7.4 (communication), \S7.1 (scalability).
Federated RL allows vehicles to share gradient updates rather than
raw sensor data, reducing bandwidth while preserving privacy.
\end{outline_note}

\subsection{Foundation Models as Universal Driving Backbones}

\begin{outline_note}
Large pre-trained models (video, language, or multi-modal) could
serve as universal feature extractors across heterogeneous traffic
scenarios, reducing the need for scenario-specific reward shaping
and enabling zero-shot transfer to new cities.
\end{outline_note}

% ==========================================================
\section{Conclusion \textnormal{\small ($\approx$600 words)}}
\label{sec:conclusion}

\begin{outline_note}
\textbf{Content.}
Synthesize the three-paradigm evolution.
Reiterate the three-layer taxonomy and how it provides a unified
framework for understanding 100 papers across the field.
Emphasize the critical bottlenecks: safety verification, scalability,
and real-time LLM integration.
End with a forward-looking statement on the convergence of MARL,
Transformer architectures, and LLM reasoning as the most promising
path toward fully autonomous, fleet-scale driving.
\end{outline_note}

% ==========================================================
\section*{References}
\addcontentsline{toc}{section}{References}

\begin{outline_note}
The bibliography will include all 100 papers from the reference list
plus foundational works cited in the body.  A \texttt{.bib} file
will be maintained alongside this document.  Citation style:
author--year (natbib), consistent with conference/journal conventions.

\textbf{Breakdown by category (planned)}:
\begin{itemize}[nosep]
  \item Foundational MARL theory: Refs 1--8 (+ Shapley 1953, Littman 1994, Bernstein 2002)
  \item Problem formulation: Refs 9--17
  \item Core MARL algorithms: Refs 18--25
  \item Communication and coordination: Refs 26--38
  \item Learning paradigms: Refs 39--50
  \item Autonomous driving applications: Refs 51--75
  \item Transformer and foundation models: Refs 76--90
  \item Safety, transfer, and scalability: Refs 83--87, 91--96
  \item Benchmarks and datasets: Refs 72--75, 82, 97--98
  \item Open challenges: Refs 99--100
  \item Papers from \texttt{refs/} folder (18 full-text papers)
\end{itemize}
\end{outline_note}

% ==========================================================
\newpage
\appendix

\section*{Appendix: Planned Figure and Table Summary}
\addcontentsline{toc}{section}{Appendix: Planned Figure and Table Summary}

\begin{table}[h!]
\centering
\caption{Summary of planned visual elements in the full survey paper.}
\label{tab:figures_summary}
\renewcommand{\arraystretch}{1.3}
\begin{tabular}{c l l l}
\toprule
\textbf{No.} & \textbf{Element} & \textbf{Location} & \textbf{Type} \\
\midrule
Fig.\ 1 & Three-paradigm evolution timeline (2003--2024) & \S1.1 & Timeline infographic \\
Fig.\ 2 & PRISMA literature selection flowchart & \S1.3 & Flowchart \\
Fig.\ 3 & Three-layer taxonomy tree & \S1.4 & Hierarchy diagram \\
Fig.\ 4 & POMDP vs.\ Dec-POMDP/POSG comparison & \S2.2 & Architecture diagram \\
Fig.\ 5 & Social LSTM vs.\ Social GAN prediction & \S3.3 & Architecture diagram \\
Fig.\ 6 & Communication topology evolution & \S4.2.3 & Flow diagram \\
Fig.\ 7 & V2X fusion strategies comparison & \S4.2.5 & System diagram \\
Fig.\ 8 & Level-$k$ game-theoretic reasoning at merge & \S4.3.3 & Scenario diagram \\
Fig.\ 9 & MAT encoder-decoder architecture & \S5.1.2 & Neural arch.\ diagram \\
Fig.\ 10 & Neuro-symbolic LLM + MPC architecture & \S5.2.4 & System diagram \\
Fig.\ 11 & Credit assignment taxonomy with trade-off curve & \S6.2 & Taxonomy + chart \\
\midrule
Table 1 & Paradigm shift summary & \S1.2 & Summary table \\
Table 2 & CTDE algorithm comparison & \S4.1.2 & Comparison table \\
Table 3 & Trajectory prediction methods & \S4.3.1 & Comparison table \\
Table 4 & Traffic signal MARL methods & \S4.4.1 & Comparison table \\
Table 5 & Simulation and benchmark ecosystem & \S4.5 & Comparison table \\
Table 6 & Safety mechanisms & \S5.3 & Comparison table \\
Table 7 & Master method comparison (all paradigms) & \S6.1 & Master table \\
\bottomrule
\end{tabular}
\end{table}

\end{document}
